% Preview source code

%% LyX 1.6.8 created this file.  For more info, see http://www.lyx.org/.
%% Do not edit unless you really know what you are doing.
\documentclass[english]{article}
\usepackage[T1]{fontenc}
\usepackage[latin9]{inputenc}
\usepackage[letterpaper]{geometry}
\geometry{verbose,tmargin=0.7in,bmargin=0.75in,lmargin=0.75in,rmargin=0.75in}
\usepackage{amsthm}
\usepackage{amsmath}
\usepackage{setspace}
\usepackage{amssymb}
\onehalfspacing

\makeatletter
%%%%%%%%%%%%%%%%%%%%%%%%%%%%%% User specified LaTeX commands.
% Preview source code

\makeatother
\date{ }

\makeatother

\usepackage{babel}

\begin{document}

%\title{Essays in International Trade Agreements}

%\maketitle

\begin{center} {\LARGE Three Essays in Applied Microeconomics}\end{center}
\begin{center} {\large Ding Liu}\end{center}

\section*{Impact of robots in monopsonistic labor markets
\vspace{-3mm}\\
{\small{}(job market paper)}}\vspace{-3mm}

Perhaps the two most important features of current labor markets in the US and other developed countries are the increasing importance of robots and the growing monopsony power of firms in labor markets. While each has been studied extensively in isolation, the interplay between them has been overlooked. This paper explores how monopsony power affects the impact of robot adoption on US labor markets. Theoretically, I show that robots have stronger effects (either positive or negative) on employment and wages in a perfectly competitive labor market than a perfectly monopsonistic labor market. Moreover, the sign of the effects depends on whether robots and labor are substitutes or complements. Empirically, I define a labor market as a commuting-zone-by-occupation cell, and measure robot exposure and labor market monpsony power for each such cell. To alleviate endogeneity concerns, I instrument for both the US exposure to robots and labor market monopsony power. My empirical results show that, from 2006 to 2014, one more industrial robot per thousand workers significantly reduces the employment-to-CZ-population ratio by 2\% and wages by 0.9\% in near-perfectly competitive labor markets. But, these employment and wage effects diminish as the labor market monopsony power grows and they become statistically insignificant in near-monopsonistic labor markets.

\section*{Bounding the joint distribution of disability and employment with misclassification
\vspace{-3mm}\\
{\small (with Daniel Millimet. \emph{Health Economics, 30(7), 1628-1647)}}}\vspace{-3mm}

Understanding the relationship between disability and employment is critical and has long been the subject of study. However, estimating this relationship is difficult, particularly with survey data, since both disability and employment status are known to be misreported. Here, we use a partial identification approach to bound the joint distribution of disability and employment status in the presence of misclassification. Allowing for a modest amount of misclassification leads to bounds on the labor market status of the disabled that are not overly informative given the relative size of the disabled population. Thus, absent further assumptions, even a modest amount of misclassification creates much uncertainty about the employment gap between the non-disabled and disabled. However, additional assumptions considered are shown to have some identifying power. For example, under our most stringent assumptions, we find that the employment gap is at least 15.2\% before the Great Recession and 22.0\% afterward.


\section*{Local labor market impacts of the 2002 Bush Steel Tariffs\vspace{-3mm}\\
{\small (with James Lake. Working paper)}}\vspace{-3mm}

President Bush announced the three-year imposition of safeguard tariffs on a variety of steel products in early 2002. Based on US local labor markets and US input-output tables, we use a difference-in-difference methodology to analyze the local labor market employment effects of these tariffs depending on how much local labor markets rely on steel as an intermediate input and how much they rely on steel production. Our results show that, at best, the tariffs only slightly boosted local employment in steel-producing industries. But, the tariffs substantially depressed local employment in steel-consuming industries and this depression did not bounce back after Bush removed the tariffs. These results suggest significant and long-lasting damage from the Trump administration's national security tariffs on steel and aluminum.

\end{document}
